\begin{result}{E.2.24}{One integral comparison for all elliptic and CM geometric factors}{Proved}
Let \(\mathcal A_{\mathrm{ell}}\) consist of all abelian
varieties over \(\Qb\) geometrically isogenous to products
of arbitrary elliptic curves and arbitrary CM abelian varieties.
Take their full integral diagram with every geometric isogeny
lattice, product, finite-model Artin module map, all Artin
objects, and the correspondence, tensor and specified projector
constructions of E.2.16.
With \(P\) and its root labels as in E.2.23, there is a
specified \(\sigma_*\in\Gamma_{F_{\mathrm{cm}}}\) with
\[
 \sigma_*|_P=(12)(34),\qquad
 \rho_A(\sigma_*)\Lambda_A=\Lambda_{{}^{\sigma_*}A}
            \quad(A\in\mathcal A_{\mathrm{ell}}).
 \tag{E.2.24a}
\]
These actual transports define one natural integral tensor
comparison on the entire stated diagram and, by E.2.3, a
comparison between the two specified real embeddings of
\(\Q(\alpha_1)\). The comparison concerns the plain integral
diagram; the labelled real Frobenius constraint of E.2.7 is
not imposed.
\end{result}
\begin{proof}
The additional issue relative to E.2.23 is that conjugation
may move an elliptic model and even its geometric isogeny
class. We first prove the precise finite-family openness
needed in this situation.

Start with any finite family of pairwise geometrically
nonisogenous non-CM elliptic curves over \(\Qb\).
Enlarge its classes to their finite \(\Gamma_\Q\)-orbits
and choose representatives \(E_i\) for the resulting finite
set. Complex conjugation permutes this set by an involution
\(\pi\). Choose actual isogenies
\(\phi_i:{}^cE_i\to E_{\pi(i)}\), and a finite Galois
field \(K/\Q\) defining all these models and isogenies.
The product representation on \(\Gamma_K\) has open image
\(G\) in the equal-determinant group \(\Delta\), by
\cite[Theorem 1.1, author PDF p.1]{GenLombardo} and the
one-factor theorem \cite[Theorem 3, p.299]{GenSerre1972}.
Set
\[
 N=\rho(\Gamma_K\cap\Gamma_{\Q^{\mathrm{cm}}}),
 \qquad S=\prod_i\operatorname{SL}_2(\widehat\Z).
\]
This is a closed normal subgroup of \(G\), contained in \(S\).

Write \(c_i\) for the actual antiholomorphic map on Betti
homology \(H_1(E_i)\to H_1({}^cE_i)\), and put
\(C_i=(\phi_i)_*c_i\). It is a rational linear isomorphism
between the indicated two homology spaces, and is integral
with nonzero determinant in their integral bases.
Naturality of the isogeny defined over \(K\) gives the exact
intertwining identity
\[
 C_i\rho_i(c\gamma c)C_i^{-1}
                 =\rho_{\pi(i)}(\gamma)
                       \quad(\gamma\in\Gamma_K).
 \tag{E.2.24b}
\]
Thus conjugation by \(c\) acts on \(G\) by
\(\theta(g)_i=C_i^{-1}g_{\pi(i)}C_i\).
Since \(K/\Q\) is Galois and \(c\) is central on
\(\Q^{\mathrm{cm}}\), the actual commutators satisfy
\[
 [\gamma,c]\in\Gamma_K\cap\Gamma_{\Q^{\mathrm{cm}}},
 \qquad g\theta(g)^{-1}\in N.
 \tag{E.2.24c}
\]
As in (E.2.23b), \(G\) contains independently supported
principal congruence subgroups at every factor and prime,
and full supported factors at all but finitely many primes.

For a two-element orbit \(\pi(i)=j\ne i\), choose such
an element \(g\) supported at \((i,\ell)\).
The element in (E.2.24c) has components \(g\) at \(i\)
and \(C_j^{-1}g^{-1}C_j\) at \(j\).
Its commutator with an independently supported \(h\) at
\((i,\ell)\) has sole component \([g,h]\) at \(i\).
Normality keeps this in \(N\). The supported-root and
derived-group identities (E.2.23e)--(E.2.23h) therefore
give an open supported subgroup at every prime, and the
full supported factor at all sufficiently large primes.
The same holds for \(j\).

For a fixed class \(\pi(i)=i\), complex uniformization
makes \(C_i\) an antilinear map \(z\mapsto a\bar z\).
Its square is positive scalar multiplication. Algebraically
that square is \(\phi_i{}^c\phi_i\in\End(E_i)=\Z\).
Comparing degrees gives
\[
 C_i^2=d_iI,\quad \operatorname{tr}C_i=0,\quad
 \det C_i=-d_i,\qquad d_i=\deg\phi_i>0.
 \tag{E.2.24d}
\]
Hence \(C_i\) is nonscalar over every \(\Q_\ell\), and
is a nonscalar invertible matrix modulo every prime
outside a finite set. In particular no local square root
of \(d_i\) is needed. Equation (E.2.24c) supplies the
supported elements \([g,C_i^{-1}]\) of \(N\).

We give both local steps explicitly. At an exceptional prime,
let \(C\in\GL_2(\Q_\ell)\) be nonscalar, and let \(R\)
be the supported subgroup of \(N\). It is normal under a
principal congruence group and contains \([g,C]\) for every
sufficiently deep supported \(g\). One of
\(X=E_{12},E_{21}\) has
\(Z=X-\operatorname{Ad}_C X\ne0\), since a matrix
commuting with both is scalar. The polynomial matrix curve
\[
 f(t)=[I+tX,C]
\]
lies in \(R\) for sufficiently small \(t\), and has
derivative \(Z\) at zero. The \(\Q_\ell\)-span of
\(\operatorname{Ad}_hZ\), with \(h\) in the normalizing
principal congruence group, is all of
\(\mathfrak{sl}_2(\Q_\ell)\). Indeed conjugation by
\(I+tE_{12}\) or \(I+tE_{21}\) is polynomial of degree
at most two in \(t\). Interpolation at three sufficiently
small distinct values makes this span stable under brackets
with both root matrices. It is a nonzero ideal, and
\([H,E_{12}]=2E_{12}\), \([H,E_{21}]=-2E_{21}\),
\([E_{12},E_{21}]=H\) show that the characteristic-zero
Lie algebra has no such proper ideal.

Choose \(h_1,h_2,h_3\) giving independent derivatives.
The polynomial map
\[
 F(t_1,t_2,t_3)=\prod_{j=1}^3h_jf(t_j)h_j^{-1}
 \tag{E.2.24e}
\]
takes a neighborhood of zero into \(R\) and has invertible
derivative. To justify the open-image step directly, use
coordinates \((a-1,b,c)\) on
\(\left(\begin{smallmatrix}a&b\\c&d\end{smallmatrix}\right)\)
near the identity, where \(d=(1+bc)/a\).
Write the coordinate map as \(Jx+P(x)\), with \(J\)
invertible and \(P\) having no constant or linear term.
On a sufficiently small complete \(\ell\)-adic ball,
\(J^{-1}P\) has Lipschitz constant less than one and
maps the ball into a smaller ball. For every \(y\) in
that ball, iteration of
\(x\mapsto y-J^{-1}P(x)\) converges and solves
\(F(x)=Jy\) in coordinates. Thus the actual image of
(E.2.24e) contains a neighborhood of the identity.
This proves openness at every exceptional prime, including
\(2\) and the nonsplit cases of (E.2.24d).

At every remaining \(\ell\geq5\), \(R\) is normal
under the full supported \(\operatorname{SL}_2(\Z_\ell)\).
Some \([g,C]\) is noncentral modulo \(\ell\).
Otherwise conjugation by \(C\) would fix each upper and
lower unipotent up to a sign. Their trace is two, which
excludes the negative sign; commuting with both roots
would make \(C\) scalar.
The required finite-group argument can be given directly.
The root commutator with \(\operatorname{diag}(2,1/2)\)
and its lower-root counterpart make
\(\operatorname{SL}_2(\mathbf F_\ell)\) perfect.
Its quotient \(\operatorname{PSL}_2(\mathbf F_\ell)\)
acts faithfully and doubly transitively on
\(\mathbf P^1(\mathbf F_\ell)\). A point stabilizer
has the normal abelian translation subgroup, whose
conjugates generate the group. Any nontrivial normal
subgroup is transitive, since its orbits are an invariant
partition for a primitive faithful action. Modulo that
subgroup every conjugate of the translation subgroup has
the same image: write an element as a product of a normal
subgroup element and a point stabilizer element.
The quotient is therefore abelian and hence trivial by
perfectness. This proves the required simplicity.
It follows that the normal reduction of \(R\) surjects
onto \(\operatorname{PSL}_2\); its remaining quotient
has order at most two and is again trivial by perfectness.
Thus \(R\) has full reduction modulo \(\ell\).

Here is the full lifting calculation. A closed subgroup
\(R\subset\operatorname{SL}_2(\Z_\ell)\) with full
reduction contains a \(g\equiv U(1)\pmod\ell\).
Write \(g=U(1)(I+\ell B)\pmod{\ell^2}\).
The first-order error in \(g^\ell\) is \(\ell\)
times the sum of \(\operatorname{Ad}_{U(-j)}B\) for
\(0\leq j<\ell\). Its entries are polynomials in
\(j\) of degree at most two, and the sums of
\(1,j,j^2\) vanish modulo \(\ell\). Hence
\[
 g^\ell\equiv I+\ell E_{12}\pmod{\ell^2}.
 \tag{E.2.24f}
\]
Conjugating by lifts in \(R\) of a lower root and an
interchange matrix gives the span of
\(E_{12},E_{21},H\) in the first congruence quotient.
For \(n\geq1\),
\[
 (I+\ell^nX)^\ell
       \equiv I+\ell^{n+1}X\pmod{\ell^{n+2}}
\]
then gives every subsequent layer. Successive approximation
and closedness prove \(R=\operatorname{SL}_2(\Z_\ell)\).

The independently supported groups from the swapped and
fixed classes therefore give an open product inside \(N\).
The closed derived subgroup of
\(\Gamma_K\cap\Gamma_{\Q^{\mathrm{cm}}}\) lies in
\(\Gamma_K\cap\Gamma_{F_{\mathrm{cm}}}\), by the
abelian base-change argument of (E.2.23g), and maps onto
the closed derived subgroup of \(N\).
The same supported perfect factors and small-prime
commutator identities used in E.2.23 make that image open.
Projection to the original finite family proves
\[
 \rho_{E_1,\ldots,E_r}
       (\Gamma_K\cap\Gamma_{F_{\mathrm{cm}}})
   \text{ is open in }
       \operatorname{SL}_2(\widehat\Z)^r.
 \tag{E.2.24g}
\]
Further finite enlargement of \(K\) preserves this:
the new compact image has finite index in the old one,
and is therefore open.

We now perform the countable construction while fixing
the changing targets. Enumerate one short Weierstrass
model \(E_n/\Qb\) from each non-CM geometric isogeny
class, using the fixed algebraic-number ordering from
E.2.17 and retaining the first model in each class.
Fix an oriented Betti basis for each model and each
conjugate by its period lattice for \(dx/y\): among all
oriented lattice bases \((\omega_1,\omega_2)\), minimize
\(|\omega_1|^2+|\omega_2|^2\), then the lexicographic
tuple of their four real coordinates. Each bounded
region contains finitely many lattice pairs, so this
specifies a unique basis. The orientation is
\(\operatorname{Im}(\bar\omega_1\omega_2)>0\).

Start with the nonempty compact quartic coset \(C_0\)
from E.2.23. Choose increasing finite Galois fields
\(K_n\) containing \(P\), the first \(n\) models,
the finite orbit enlargements and isogenies used in
(E.2.24g), and the \(n\)-th member of the fixed finite
Galois exhaustion of \(\Qb\).
These choices can be fixed by the first appropriate
field and isogeny in the countable algebraic-data
enumerations. Partition the preceding compact fibre
by restriction to \(K_n\), and select the first
nonempty finite cylinder \(C'_n\).
The models \({}^{\sigma}E_i\) and all target bases
are now fixed for \(\sigma\in C'_n\).

For any \(\sigma_n\in C'_n\), its stabilizer under
right multiplication is exactly
\[
 H_n=\Gamma_{F_{\mathrm{cm}}}\cap\Gamma_{K_n}
                  \cap\bigcap_{i<n}\ker\rho_{E_i}.
 \tag{E.2.24h}
\]
Indeed, for \(k\in\Gamma_{K_n}\), actual transport
satisfies
\(\rho_{E_i}(\sigma_nk)=
\rho_{E_i}(\sigma_n)\rho_{E_i}(k)\).
Thus equality of all previous matrices is precisely
the kernel condition in (E.2.24h).
Openness (E.2.24g) makes \(\rho_{E_n}(H_n)\) open,
because its product congruence neighborhood contains
each sufficiently deep last-factor element with the
previous components exactly the identity.
The possible new matrices form the open coset
\[
 \rho_{E_n}(\sigma_n)\rho_{E_n}(H_n)
                \subset\operatorname{SL}_2(\widehat\Z).
 \tag{E.2.24i}
\]
The determinant is one since \(\sigma_n\) fixes all
roots of unity and the bases are oriented.
By the proved elementary density in E.2.23 this coset
contains an integral \(\operatorname{SL}_2(\Z)\)
matrix. Choose the first in entry-height and
lexicographic order, and let \(C_n\subset C'_n\)
be its inverse image. These are nested nonempty
compact sets. Their intersection contains an element
\(\sigma_*\), and contains only one: the selected
\(K_n\)-restrictions exhaust \(\Qb\).
It retains the prescribed quartic action and has
integral actual transport on every \(E_n\).
This specifies an element without claiming an
effective decision procedure for the choices.

On products of these elliptic representatives and
CM factors, take the direct sum of their actual
integral transports and the idele-one CM comparison.
For every geometric isogeny \(q:B\to A\), the exact
superlattice (E.2.23l) and its finite-kernel
interpretation show that
\(({}^{\sigma_*}q)_*T_Bq_*^{-1}\) carries
\(\Lambda_A\) onto \(\Lambda_{{}^{\sigma_*}A}\).
Completion injectivity gives independence of \(q\)
and every correspondence, tensor and projector
equation. Inverse duals and exterior powers give
cohomology, and the Tate object is fixed.
All finite-model Artin blocks and their module maps
are compatible by E.2.3 and E.2.16. This proves
(E.2.24a) on the entire specified diagram.

The theorem supplies no integrality for additional
simple non-CM abelian factors of dimension at least
two and no extension to every pure motive.
Those are still part of the full original
place-separation problem.
\end{proof}
